\documentclass{article}
\usepackage{amsmath,amssymb}
\usepackage{xurl}
\usepackage[hidelinks]{hyperref}
\setlength{\emergencystretch}{2em}
% Shared commands for notes in docs/paper-gaps/.

\DeclareUnicodeCharacter{2080}{\ifmmode{}_{0}\else\textsubscript{0}\fi}
\DeclareUnicodeCharacter{2081}{\ifmmode{}_{1}\else\textsubscript{1}\fi}
\DeclareUnicodeCharacter{2082}{\ifmmode{}_{2}\else\textsubscript{2}\fi}
\DeclareUnicodeCharacter{2099}{\ifmmode{}_{n}\else\textsubscript{n}\fi}

\newcommand{\Lean}{\textsc{Lean}}
\ExplSyntaxOn
\NewDocumentCommand{\leanid}{m}{
  \ifmmode
    \textsf{\detokenize{#1}}
  \else
    \group_begin:
    \urlstyle{sf}
    \tl_set:Nn \l_tmpa_tl {#1}
    \tl_replace_all:Nnn \l_tmpa_tl {\_} {_}
    \exp_args:NV \path \l_tmpa_tl
    \group_end:
  \fi
}
\ExplSyntaxOff
\providecommand{\tr}{\operatorname{tr}}
\newcommand{\tnleanIssueBase}{https://github.com/LionSR/TNLean/issues/}
\newcommand{\tnleanPrBase}{https://github.com/LionSR/TNLean/pull/}
\newcommand{\tnleanDocBase}{https://sirui-lu.com/TNLean/docs/}
\newcommand{\ghissue}[1]{\href{\tnleanIssueBase#1}{GitHub issue \##1}}
\newcommand{\ghpr}[1]{\href{\tnleanPrBase#1}{PR \##1}}
\newcommand{\doclink}[1]{\href{\tnleanDocBase#1.html}{\path{#1}}}

% Dirac notation and the identity operator, as in blueprint/src/macros/common.tex.
% \providecommand keeps note-local definitions authoritative.
\usepackage{dsfont}
\providecommand{\ket}[1]{|#1\rangle}
\providecommand{\bra}[1]{\langle #1|}
\providecommand{\braket}[2]{\langle #1 | #2 \rangle}
\providecommand{\ketbra}[2]{|#1\rangle\!\langle #2|}
\providecommand{\Id}{\mathds{1}}
\providecommand{\one}{\mathds{1}}
\providecommand{\C}{\mathbb{C}}
\providecommand{\MN}[1]{M_{#1}(\C)}
\providecommand{\MD}{M_D(\C)}

% External referencing for paper sources.  The build helper writes sanitized
% auxiliary files under build/paper-aux/ and excludes duplicated source labels.
\usepackage{xr}

% Import a generated paper auxiliary file with a caller-supplied label prefix.
% The candidate paths support builds from docs/paper-gaps/, the repository root,
% and build/paper-aux/ itself.
\newcommand{\importpaperaux}[2]{%
  \IfFileExists{../../build/paper-aux/#2.aux}{%
    \externaldocument[#1]{../../build/paper-aux/#2}%
  }{%
    \IfFileExists{build/paper-aux/#2.aux}{%
      \externaldocument[#1]{build/paper-aux/#2}%
    }{%
      \IfFileExists{#2.aux}{%
        \externaldocument[#1]{#2}%
      }{}%
    }%
  }%
}

% General external reference macros
\newcommand{\paperref}[2]{\ref{#1-#2}}
\newcommand{\papereqref}[2]{(\ref{#1-#2})}

% CPSV16 source (1606.00608 / MPDO-22-12-17-2.tex) external document and shortcuts
\importpaperaux{cpsv16-}{1606.00608/MPDO-22-12-17-2-tnlean}
\newcommand{\cpsvref}[1]{\paperref{cpsv16}{#1}}
\newcommand{\cpsveqref}[1]{\papereqref{cpsv16}{#1}}

% Verdict marker: \gapnote{<kind>}{<status>}. Kind is the policy
% classification (clarification, local-correction, scope-restriction,
% unfaithful, false-source, open-gap); status is open, wip (elimination
% underway), resolved, or historical. Machine-read by the site index;
% prints a verdict line.
\newcommand{\gapnote}[2]{\par\noindent\textbf{Verdict:} #1 (#2).\par}

\title{Matrix Product Operator Symmetric States: the Periodic-Boundary Layer}
\author{}
\date{2026-09-24}
\begin{document}
\maketitle
\gapnote{scope-restriction}{open}

\section{The source definitions}
Garre-Rubio, Lootens and Moln\'ar define a matrix product operator
algebra through the closedness of the arbitrary-boundary operators
$\mathcal A_T^n=\{O^n_{T,B}: B\in\mathcal M_\chi\}$ under multiplication,
and a symmetric matrix product state through the invariance
$\mathcal A_T\cdot\mathcal S_A\subset\mathcal S_A$ of the arbitrary-boundary
subspace $\mathcal S_A^n=\{\lvert\psi^n_{A,X}\rangle: X\in\mathcal M_D\}$ for every
system size~$n$~\cite{GarreRubio2023MPOSym}.\footnote{Source traceability:
    \path{Papers/2203.12563/REsubmission.tex}, lines 311--345 (the subspace
    $\mathcal S_A^n$ and the closedness condition) and lines 431--434
    (the definition of a symmetric state).}
The periodic operators $O_a=O_{T_a,\mathbb 1}$ of the blocks $a$ then
satisfy
\begin{align}
    O_aO_b & =\sum_c N_{ab}^c O_c
    \label{eq:glm23_fusion}
\end{align}
for every system size, with $N_{ab}^c\in\mathbb N$, and the action of a block
$a$ on a block $x$ of the state decomposes into blocks $y$ with multiplicities
$M_{a,x}^y\in\mathbb N$ through the action tensors. Associativity of the action
gives
\begin{align}
    \sum_c N_{ab}^c M_{c,x}^y & =\sum_z M_{b,x}^z M_{a,z}^y .
    \label{eq:glm23_nimrep}
\end{align}
\footnote{Source traceability: lines 361--362 for
    (\ref{eq:glm23_fusion}); lines 458--460 for the multiplicities; lines
    491--492 for associativity; line 683 for the group case
    $M_{gh,x}^y=\sum_z M_{g,z}^y M_{h,x}^z$.}
For groups the source also formulates a weaker, periodic-boundary version
of both conditions, in which the boundaries are proportional to the identity
in each block.\footnote{Source traceability: lines 994--1000 and 1062--1064.}

\section{What is formalized}
The formal layer records only periodic-boundary data, for arbitrary finite
fusion rules. The fusion hypothesis is (\ref{eq:glm23_fusion}) at every
positive length. The symmetry hypothesis on a family of tensors $A_x$ is
\begin{align}
    O_a\lvert\psi_{A_x}\rangle & =\sum_y M_{a,x}^y\lvert\psi_{A_y}\rangle
    \qquad\text{at every positive length,}
    \label{eq:glm23_periodic_symmetry}
\end{align}
with complex coefficients independent of the length. Taking identity
boundaries in the action-tensor decomposition shows that the source's
definition implies (\ref{eq:glm23_periodic_symmetry}), so the formal
hypothesis is weaker than the source's. From it the formal results derive, for
normal blocks, that each $M_{a,x}^y$ is a nonnegative integer and that
(\ref{eq:glm23_nimrep}) holds; the source takes integrality from the action
tensors instead.\footnote{Formal definitions:
    \leanid{MPOTensor.IsMPOFusionAlgebra},
    \leanid{MPOTensor.IsMPOSymmetricFamily}, \leanid{MPOTensor.IsMPOSymmetric}.
    Results: \leanid{MPOTensor.exists_isNIMRep_of_isMPOSymmetricFamily} and
    \leanid{MPOTensor.exists_mpo_mulVec_eq_of_isInvertibleLabel}.}

The source's standing assumption that the blocks are injective and
orthogonal to each other in the thermodynamic limit (line 317) is used in
the form that the periodic vectors of the blocks are linearly independent at
one positive length. Asymptotic orthogonality implies this at every
sufficiently large length. It is the same reading as the eventual linear
independence clause of the formal basis of normal tensors.

\section{Boundary}
Not formalized: the arbitrary-boundary algebra $\mathcal A_T^n$, the
subspace $\mathcal S_A^n$ and the invariance $\mathcal A_T\cdot\mathcal
S_A\subset\mathcal S_A$; the fusion and action tensors; the $F$- and
$L$-symbols and the pentagon equations. Consequently the single-block
obstruction of the source's subsection ``Injective MPSs cannot be invariant
under non-trivial MPOs'', which shows that the $F$-symbols
are trivial for an injective invariant state without
multiplicity,\footnote{Source traceability: lines 607--636 and line 738.}
is formalized only in its fusion-ring form: a single normal symmetric block on
which the unit acts trivially gives a ring homomorphism from the fusion ring to
$\mathbb Z$ that is nonnegative on the labels. This excludes the Fibonacci
algebra, but it does not exclude group-like algebras with a nontrivial
three-cocycle.

Elimination plan: define the arbitrary-boundary algebra and subspace,
prove that the source's invariance implies
(\ref{eq:glm23_periodic_symmetry}), and state the cohomological single-block
obstruction through the action tensors of the multi-block compression
theorem.
\bibliographystyle{alpha}
\bibliography{references}
\end{document}
